\documentclass[12pt]{article}

\setlength{\parindent}{0pt}
\setcounter{secnumdepth}{3}
\usepackage[margin=1in]{geometry}

% packages
\usepackage{amsmath, amsfonts, amsthm, amssymb}
\usepackage{graphicx, float}
\usepackage{mathtools}
\usepackage{titlesec}
\usepackage{interval}
\usepackage{hyperref}
\usepackage{siunitx}
\usepackage{titling}
\usepackage{vwcol}
\usepackage{setspace}
\usepackage{empheq}
\usepackage{cancel}
\usepackage{esdiff}
\usepackage{multicol}
\usepackage{mdframed}
\usepackage{esdiff}
\usepackage{tikzsymbols}
\usepackage{multicol}
\usepackage{tikz}
\usepackage{varwidth}
\usepackage{pgfplots}
\usepackage{fancyhdr}
\usepackage{enumitem}
\usepackage{tcolorbox}

\intervalconfig {
	soft open fences
}

% header
\pagestyle{fancy}
\fancyhf{}
\lhead{Ethan Fan}
\rhead{AP Calculus BC}
\cfoot{\thepage}

% section format
\titleformat{\section}
  {\large \bfseries}
  {}
  {0em}
  {}
\titlespacing*{\section}{0pt}{0.5em}{0.3em}

\titleformat{\subsection}[runin]
  {\bfseries}
  {}
  {0em}
  {}
\titlespacing*{\subsection}{0pt}{0pt}{0em}

\titleformat{\subsubsection}[runin]
  {\itshape}
  {(\roman{subsubsection})}
  {0em}
  {}

% commands
\newcommand{\alignedintertext}[1]{%
  \noalign{%
    \vskip\belowdisplayshortskip
    \vtop{\hsize=\linewidth#1\par
    \expandafter}%
    \expandafter\prevdepth\the\prevdepth
  }%
}

\newcommand*{\paren}[1]{\ensuremath\left(#1\right)}
\newcommand*{\limit}[2][x]{\ensuremath{\displaystyle\lim_{#1\to#2}}}
\newcommand*{\Limit}[3][x]{\ensuremath{\displaystyle\lim_{#1\to#2}\left[#3\right]}}
\newcommand*{\deriv}[1][x]{\ensuremath{\dfrac{\mathrm{d}}{\mathrm{d}#1}}}
\newcommand*{\Deriv}[2][x]{\ensuremath{\dfrac{\mathrm{d}}{\mathrm{d}#1}\left[#2\right]}}
\newcommand{\ddx}{\dfrac{d}{dx}}
\newcommand{\dydx}{\dfrac{dy}{dx}}
\newcommand{\dydt}{\dfrac{dy}{dt}}
\newcommand{\dxdt}{\dfrac{dx}{dt}}
\newcommand*{\iinteg}[2][x]{\ensuremath{\displaystyle\int #2\;\mathrm{d}#1}}
\newcommand*{\dinteg}[4][x]{\ensuremath{\displaystyle\int_{#2}^{#3}#4\;\mathrm{d}#1}}
\newcommand*{\abs}[1]{\ensuremath{\left|#1\right|}}
\newcommand*{\eps}{\varepsilon}
\newcommand*{\real}{\ensuremath\mathbb{R}}
\newcommand*{\integer}{\ensuremath\mathbb{Z}}
\newcommand*{\posint}{\ensuremath\mathbb{Z^+}}
\newcommand*{\cbrt}[1]{\ensuremath{\sqrt[3]{#1}}}
\newcommand*{\lheqzero}{\ensuremath{\underset{\text{L'H}}{\overset{\left[\frac00\right]}{=}}}}
\newcommand*{\lheqinfty}{\ensuremath{\underset{\text{L'H}}{\overset{\left[\frac{\infty}{\infty}\right]}{=}}}}
\newcommand{\tor}{\text{ or }}
\newcommand{\floor}[1]{\lfloor #1 \rfloor}

\newcommand{\vem}{\vspace{0.5em}}
\newcommand{\example}[1]{\section{Example #1}}
\newcommand{\problem}[1]{\section{Problem #1}}
\newcommand{\subproblem}[1]{\subsection{(#1)} \,}

\DeclareMathOperator{\dne}{DNE}
\DeclareMathOperator{\sgn}{sgn}

\newcommand{\definition}[1]{\begin{tcolorbox}[colback=red!5!white,colframe=red!75!black,parbox=false] #1 \end{tcolorbox}}

\newcommand{\theorem}[2]{\begin{tcolorbox}[title={#1},colback=blue!5!white,colframe=blue!75!black,parbox=false] #2 \end{tcolorbox}}

\allowdisplaybreaks
% \postdisplaypenalty=100000

% document
\begin{document}

\vspace*{0cm}

% opening
\begin{center}
    {\Large \bfseries Problem Set 7} \\[0.5cm]
    {\large The Net Change Theorem} \\[0.35cm]
    {\normalsize September 18, 2026}
\end{center}

% problems

\problem{5}

\subproblem{a}
\begin{gather*}
    \boxed{v(t)=t^2+3t-4}
\end{gather*}

\subproblem{b}
\begin{gather*}
    v(t)=(t-1)(t+4)\\
    x(t)=\frac{1}{3}t^3+\frac{3}{2}t^2-4t+C\\
    \int_0^3 |v(t)|\, dt = \int_1^3 v(t)\, dt- \int_0^1 v(t)\, dt\\
    =x(3)-2\cdot x(1) \\
    =9+13.5-12-2(\frac{1}{3}+\frac{3}{2}-4)\\
    =\boxed{\frac{89}{6}m}
\end{gather*}

\problem{7}

\subproblem{a}
\begin{gather*}
    \max: 230\\
    \min: 172
\end{gather*}
The increasing condition is required, as the estimates are dependent on the left and right ends of the interval being the maximums and the minimums. \vem

\subproblem{b}
\begin{gather*}
    Q(6)=6+17+30+41+50+57=201
\end{gather*}

\problem{8}

\subproblem{a}
\begin{gather*}
    A(t)=200+\int_0^t R(t)-D(t) \, dt\\
    =200+\int_0^t40-t-\frac{t^2}{20}\, dt\\
    =200+(-\frac{t^3}{60} - \frac{t^2}{2}+40t) \Big|^t_0\\
    A(t)=200+40t-\frac{t^2}{2}-\frac{t^3}{60}
\end{gather*}

\subproblem{b}
\begin{gather*}
    A'(t)=40-t-\frac{t^2}{20}\\
    A'(10)>0\\
    A'(25)<0
\end{gather*}
The amount of water is increasing at $t=10$, but is decreasing at $t=25$.\vem

\subproblem{c}
\begin{gather*}
    A'(t)=0 \implies 40-t-\frac{t^2}{20}=0\\
    t^2+20t-800=0\\
    (t+40)(t-20)=0 \implies t=20\\
    A(20)=200+800-200-\frac{400}{3}=\frac{2000}{3}
\end{gather*}
The amount of water is greatest at $t=20$, with $\dfrac{2000}{3}$ liters of water.\vem

\subproblem{d}
\begin{gather*}
    t=30 \implies A(30)=500
\end{gather*}
The net change is not the same as the total amount of water entering, as the net change considers both the amount of water entering and leaving.

\problem{9}

\subproblem{a}
\begin{gather*}
    P'(x)=3-0.01x+0.000006x^2\\
    P(x)=3x-0.005x^2+0.000002x^2\\
    \int_{2000}^{4000}P'(x)\, dx\\
    =P(4000)-P(2000)\\
    =60000-2000\\
    =\boxed{58000}
\end{gather*}

\subproblem{b}
\begin{gather*}
    P(t)=300\sin \paren{\frac{\pi t}{12}}\\
    \int_0^{12} P(t)\, dt\\
    =-\frac{3600}{\pi}\cos \paren{\frac{\pi t}{12}} \Big|^{12}_0\\
    =\frac{7200}{\pi}\text{ Wh}\\
    =\boxed{\frac{7.2}{\pi} \text{ kWh}}
\end{gather*}

\problem{10}
\begin{gather*}
    3\int_0^a e^x\, dx= \int_0^b e^x\, dx\\
    3e^a-3=e^b-1\\
    3e^a-2=e^b\\
    \boxed{b=\ln (3e^a-2)}
\end{gather*}

\problem{11}

\subproblem{a}
\begin{gather*}
    D(T)=\int_0^T v(t)\, dt=\int_0^T t^2-4t+3\, dt= \frac{T^3}{3}-2T^2+3T\\
    D(T)=\frac{T}{3}(T-3)^2\\
    D(T)=0 \implies \boxed{T=3}
\end{gather*}

\subproblem{b}
\begin{gather*}
    T\ge 0 \implies D(T)=\frac{T}{3}(T-3)^2\ge 0
\end{gather*}
The particle is always at or to the right of its starting position. \vem

\subproblem{c}
\begin{gather*}
    \int_0^3|v(t)|\, dt=\int_0^1v(t)\, dt-\int_1^3 v(t)\, dt = 2\cdot D(1)-D(3)= \boxed{\frac{8}{3}}
\end{gather*}
The total distance traveled is 8 meters. The particle can travel for the same distance in opposite directions to attain zero displacement.\vem

\subproblem{d}
\begin{gather*}
    \frac{8}{3}+\int_3^T v(t)\, dt=4\\
    \frac{T}{3}(T-3)^2=\frac{4}{3}\\
    \boxed{T=4}
\end{gather*}

\subproblem{e}
\begin{gather*}
    v(t)=(t-a)(t-b)=t^2-(a+b)t+ab\\
    D(T)=\int_0^T v(t)\, dt=\int_0^T t^2-(a+b)t+ab \, dt= \frac{T^3}{3}-\frac{(a+b)T^2}2+abT\\
    D(b)=\frac{b^3}{3}-\frac{ab^2+b^3}{2}+ab^2=\frac{b^2}{6}(3a-b)\\
    D(b)=0 \implies \boxed{b=3a}\\
    \int_0^{3a} |v(t)|\, dt = \int_0^a v(t)\, dt- \int_a^{3a} v(t)\, dt\\
    =2\cdot D(a)-D(3a)=2\paren{\frac{a^3}{3}-2a^3+3a^3}=\boxed{\frac{8a^3}{3}}
\end{gather*}










\end{document}

